% Generated draft: requires mathematical and source review.
\documentclass[11pt]{article}
\usepackage[margin=25mm]{geometry}
\usepackage{fontspec}
\setmainfont{Latin Modern Roman}
\usepackage{amsmath,amssymb,mathtools,unicode-math}
\setmathfont{Latin Modern Math}
\usepackage{enumitem,needspace,xurl,hyperref}
\setlength{\parindent}{0pt}
\setlength{\parskip}{5pt}
\setlength{\emergencystretch}{3em}
\newcommand{\N}{\mathbb{N}}
\newcommand{\Z}{\mathbb{Z}}
\newcommand{\Q}{\mathbb{Q}}
\newcommand{\R}{\mathbb{R}}
\newcommand{\C}{\mathbb{C}}
\newcommand{\F}{\mathbb{F}}
\newcommand{\A}{\mathbb{A}}
\newcommand{\PP}{\mathbb{P}}
\newcommand{\E}{\mathbb{E}}
\newcommand{\eps}{\varepsilon}
\newcommand{\dd}{\,\mathrm{d}}
\begin{document}
\section*{UnsolvedMath ID: 3000050}
\textbf{Source problem code:} catalogue entry
\subsection{Definitions and mathematical statement}
Let \(D=(V,E)\) be a finite digraph whose arc set can be partitioned into two directed spanning trees. ``Disjoint'' here must mean arc-disjoint, since both trees span \(V\). The source asks for a characterization of when there exists a possibly different partition
\[
E=T\mathbin{\dot\cup}(E\setminus T)
\]
into directed spanning trees such that reversing precisely the arcs of \(T\) gives a strongly connected digraph \(D_T\). Reversal preserves each arc's identity and exchanges its head and tail; all other arcs retain their orientations. Strong connectivity means a directed path exists from every vertex to every other vertex.

For \(S\subset V\), let \(\delta_D^+(S)\) and \(\delta_D^-(S)\) be arcs leaving and entering \(S\). The requested feasibility property can be stated explicitly as existence of a tree partition satisfying
\[
|\delta_D^+(S)\setminus T|+|\delta_D^-(S)\cap T|\ge1
\qquad(\varnothing\ne S\subsetneq V),
\]
which is exactly the no-empty-directed-cut condition for \(D_T\). This merely specifies the property to characterize; it is not offered as the sought structural characterization or algorithm.

A directed spanning tree may mean an arbitrary orientation of an undirected spanning tree, or a rooted in-/out-arborescence with orientation constraints. The excerpt gives no definition, roots, or direction convention. Under the former reading each underlying edge set is a spanning tree; under an arborescence reading additional indegree/outdegree and root constraints must be retained. Loops and parallel-arc conventions are also unspecified. No polynomial-time or min--max characterization requirement is supplied. The research summary reports OPEN-TRIAGE, saying no correct characterization or resolving citation was verified.
\subsection{Short English statement}
Characterize when a digraph decomposable into two spanning trees has such a decomposition for which reversing one tree makes the whole digraph strongly connected. The meaning of directed tree and the intended characterization format are unspecified.
\subsection{Sources}
\begin{itemize}
\item Local UnsolvedMath registry, permanent ID \(3000050\), exact source code \texttt{catalogue entry}, ``Making the union of two directed spanning trees strongly connected.''
\item Egres Open, \emph{List of open problems}, page of the same title; supplied index \url{https://oldlemon.cs.elte.hu/egres/open/List_of_open_problems}. HTML extraction; recorded access 2026-07-29.
\item No authoritative research-paper citation or definition of the directed-tree convention is supplied; the research summary reports only an unverified no-resolution search.
\end{itemize}
\end{document}
